% !TEX root = ../Presentation.tex
% ---------------------------------------------------------------------------
% The three models and how they are checked
% ---------------------------------------------------------------------------
\section{Models}

\begin{frame}{Three models}
  \begin{itemize}
    \setlength{\itemsep}{1em}
    \item \textbf{Model A:} one curve per group, all sessions pooled
    \item \textbf{Model B:} a curve for every session\\
          {\color{muted}sessions may differ, but are assumed to resemble each other}
    \item \textbf{Model C:} like B, but the distractor's shift is proportional to its
          amplitude\\
          {\color{muted}one curve per hand, with both trial types}
  \end{itemize}
\end{frame}

\begin{frame}{How I check a model}
  \begin{enumerate}
    \setlength{\itemsep}{1em}
    \item \textbf{Simulated data:} make data with known answers, and check that the model
          gets them back
    \item \textbf{Real data:} check that the fitted model reproduces what the monkey did
    \item \textbf{Compare models:} how well each predicts trials left out of its fit
  \end{enumerate}
\end{frame}
